\documentclass[11pt]{article}

\usepackage[margin=1in]{geometry}
\usepackage{xcolor}
\usepackage{listings}
\usepackage{hyperref}
\usepackage{booktabs}
\usepackage{enumitem}
\usepackage{graphicx}

\definecolor{codebg}{RGB}{245,245,245}
\definecolor{codekw}{RGB}{0,0,180}
\definecolor{codecm}{RGB}{96,96,96}
\definecolor{codestr}{RGB}{163,21,21}

\lstdefinestyle{cppstyle}{
  language=C++,
  basicstyle=\ttfamily\footnotesize,
  backgroundcolor=\color{codebg},
  keywordstyle=\color{codekw}\bfseries,
  commentstyle=\color{codecm}\itshape,
  stringstyle=\color{codestr},
  numbers=left,
  numberstyle=\tiny\color{codecm},
  numbersep=8pt,
  showstringspaces=false,
  breaklines=true,
  frame=single,
  framerule=0.4pt,
  tabsize=2,
  captionpos=b,
  xleftmargin=12pt,
  morekeywords={Request,ReplicaInfo,NetChannel,ResDBConfig,std,uint64_t,size_t,absl}
}
\lstdefinestyle{shstyle}{
  basicstyle=\ttfamily\footnotesize,
  backgroundcolor=\color{codebg},
  numbers=none,
  showstringspaces=false,
  breaklines=true,
  frame=single,
  framerule=0.4pt,
  xleftmargin=12pt,
  captionpos=b,
}

\title{\textbf{CS 410 -- Large Scale Systems}\\Assignment 3: Implementing MiniSpanner\\Sharded--Replicated Architecture with Two--Phase Commit}
\author{Dom Ellis}
\date{\today}

\begin{document}
\maketitle

\section{Overview}

This report describes our implementation of a sharded--replicated architecture
on top of Apache ResilientDB (incubating), augmented with the Two--Phase
Commit (2PC) protocol required for the CS~410 track. The system is structured
as four shards of four replicas each, for a total of sixteen
\texttt{kv\_service} processes. The four shard leaders participate in 2PC to
decide the fate of every distributed transaction; each shard then runs PBFT
internally to replicate the committed write across its four members.

The deliverable code lives on the \texttt{DomTest} branch of our fork of
ResilientDB. The runtime artifact is a Docker image, \texttt{dom-sharded},
whose default \texttt{ENTRYPOINT} brings the entire 16--node cluster up on
\texttt{127.0.0.1} ports \texttt{18001--18034}.

\section{Architecture}

\subsection{Topology}
\begin{itemize}[itemsep=2pt]
  \item \textbf{Shards}: 4 (region IDs 1--4).
  \item \textbf{Replicas per shard}: 4 (16 \texttt{kv\_service} processes total).
  \item \textbf{Port layout}: shard $s$ replica $r$ listens on
        \texttt{180\{s\}\{r\}} (e.g., shard~1 uses 18001--18004; shard~2 uses
        18011--18014; shard~3 uses 18021--18024; shard~4 uses 18031--18034).
  \item \textbf{Shard leaders} (primaries): node IDs \texttt{1, 5, 9, 13}.
\end{itemize}

\subsection{End-to-end transaction path}

\begin{lstlisting}[style=shstyle,caption={High-level flow of a single SET transaction.}]
kv_service_tools (round-robin across 4 leaders)
   ---> shard leader receives TYPE_CLIENT_REQUEST
        ---> ResponseManager::NewUserRequest()  enqueues in batch_queue_
             ---> BatchProposeMsg() drains, forwards local primary as TYPE_NEW_TXNS
                  ---> Commitment::ProcessNewRequest()
                       [2PC] PREPARE  ---direct NetChannel---> 3 cross-shard leaders
                                            ---> Process2PCPrepare ---> VOTE YES
                       [2PC] collect 3/3 votes
                       [2PC] GLOBAL_COMMIT ---> 3 cross-shard leaders
                                            ---> Process2PCCommit ---> start local PBFT
                       PBFT on coordinator shard
                       ---> PRE_PREPARE -> PREPARE -> COMMIT -> execute -> client reply
\end{lstlisting}

\subsection{Configuration files}

The deployment uses static JSON configuration committed at
\texttt{scripts/deploy/config/sharded/}:

\begin{itemize}[itemsep=2pt]
  \item \texttt{server/shard\{1..4\}.server.config} -- one per shard. Each file
        lists the four local replicas under \texttt{region} and declares the
        other three shard leaders under a new \texttt{crossShardPeer} field.
  \item \texttt{client.shard\_leaders.config} -- used by the
        \texttt{kv\_service\_tools} client. Contains the four shard leaders
        (one per region) and sets
        \texttt{multiShardClientRoundRobin: true}.
\end{itemize}

\section{Changes to ResilientDB}

\subsection{Schema additions (Protocol Buffers)}

Three new message types and a cross-shard peer list were added to the
existing protos. The dispatcher learns the new types from
\texttt{platform/proto/resdb.proto}; the configuration loader learns the
cross-shard topology from \texttt{platform/proto/replica\_info.proto}.

\begin{lstlisting}[style=cppstyle,caption={\texttt{platform/proto/resdb.proto} -- new request types added to \texttt{Request::Type}.}]
TYPE_2PC_PREPARE = 22;
TYPE_2PC_VOTE    = 23;
TYPE_2PC_COMMIT  = 24;
NUM_OF_TYPE      = 25;  // bumped from 22 to 25
\end{lstlisting}

\begin{lstlisting}[style=cppstyle,caption={\texttt{platform/proto/replica\_info.proto} -- cross-shard topology and client round-robin flag added to \texttt{ResConfigData}.}]
message ResConfigData {
  // ... existing fields ...
  optional bool multi_shard_client_round_robin = 25;
  repeated ReplicaInfo cross_shard_peer = 26;
}
\end{lstlisting}

\subsection{Configuration accessors}

Two helpers were added to \texttt{ResDBConfig} so the runtime can read the
new JSON fields.

\begin{lstlisting}[style=cppstyle,caption={\texttt{platform/config/resdb\_config.cpp} -- new accessors for the 2PC topology and round-robin policy.}]
bool ResDBConfig::MultiShardClientRoundRobin() const {
  return config_data_.has_multi_shard_client_round_robin() &&
         config_data_.multi_shard_client_round_robin();
}

std::vector<ReplicaInfo> ResDBConfig::GetCrossShardPeers() const {
  std::vector<ReplicaInfo> out;
  out.reserve(static_cast<size_t>(config_data_.cross_shard_peer_size()));
  for (const auto& p : config_data_.cross_shard_peer()) {
    out.push_back(p);
  }
  return out;
}
\end{lstlisting}

\subsection{Client-side round-robin across shard leaders}

The client's \texttt{PickDestReplica()} in the transaction constructor
alternates across the shard leaders listed in the client config when
the \texttt{multiShardClientRoundRobin} flag is set. A static atomic
counter advances on each request, modulo the number of leaders, so a
single long-lived client cycles
$1 \rightarrow 2 \rightarrow 3 \rightarrow 4 \rightarrow 1 \rightarrow \ldots$
through the four shards, exactly as required by the assignment. The
counter is \texttt{static} so all \texttt{KVClient} instances in one
process share a single rotor.

\begin{lstlisting}[style=cppstyle,caption={\texttt{interface/rdbc/transaction\_constructor.cpp} -- proxy round-robin across shard leaders.}]
void TransactionConstructor::PickDestReplica() {
  const std::vector<ReplicaInfo>& replicas = config_.GetReplicaInfos();
  if (replicas.empty()) {
    return;
  }
  if (config_.MultiShardClientRoundRobin() && replicas.size() > 1) {
    const uint64_t round =
        proxy_send_round_.fetch_add(1, std::memory_order_relaxed);
    const size_t idx = static_cast<size_t>(round % replicas.size());
    NetChannel::SetDestReplicaInfo(replicas[idx]);
  } else {
    NetChannel::SetDestReplicaInfo(replicas[0]);
  }
}

int TransactionConstructor::SendRequest(
    const google::protobuf::Message& message, Request::Type type) {
  PickDestReplica();
  return NetChannel::SendRequest(message, type, false);
}
\end{lstlisting}

\subsection{Dispatching the new 2PC messages}

\texttt{ConsensusManagerPBFT::InternalConsensusCommit} is the central
\texttt{switch} that routes every inbound \texttt{Request} by type.
Three arms route the 2PC message types to the corresponding handlers in
\texttt{Commitment}.

\begin{lstlisting}[style=cppstyle,caption={\texttt{platform/consensus/ordering/pbft/consensus\_manager\_pbft.cpp} -- new dispatcher arms for the 2PC types.}]
case Request::TYPE_2PC_PREPARE:
  return commitment_->Process2PCPrepare(std::move(context),
                                        std::move(request));
case Request::TYPE_2PC_VOTE:
  return commitment_->Process2PCVote(std::move(context),
                                     std::move(request));
case Request::TYPE_2PC_COMMIT:
  return commitment_->Process2PCCommit(std::move(context),
                                       std::move(request));
\end{lstlisting}

\subsection{Coordinator: initiating 2PC}

When a shard leader's \texttt{Commitment::ProcessNewRequest} receives a new
batch from \texttt{ResponseManager}, it now broadcasts a \texttt{2PC PREPARE}
to its three cross-shard peers \emph{before} starting local PBFT. The
PREPARE carries the coordinator's contact information in the
\texttt{client\_info} field so the participants know where to send their
vote back.

\paragraph{Cross-shard transport.} The standard intra-shard
\texttt{ReplicaCommunicator} pool is designed for the same-shard peer set
established at startup -- it remaps destinations to a long-connection
port and signs payloads with the local shard's verifier. Cross-shard
peers live on the base listener port and have no shared public keys, so
we use a fresh short-connection \texttt{NetChannel(ip, port)} for the
three cross-shard sends (PREPARE, VOTE, GLOBAL\_COMMIT). Intra-shard
PBFT continues to use the existing \texttt{ReplicaCommunicator} pool.

\begin{lstlisting}[style=cppstyle,caption={\texttt{platform/consensus/ordering/pbft/commitment.cpp} -- coordinator broadcasts 2PC PREPARE then stashes the request awaiting votes.}]
// === 2PC Phase 1: PREPARE to participants ===
LOG(ERROR) << "[2PC] Coordinator starting 2PC for seq: " << *seq;
const std::vector<ReplicaInfo> cross_peers = config_.GetCrossShardPeers();
ReplicaInfo coord_contact;
coord_contact.set_id(config_.GetSelfInfo().id());
coord_contact.set_ip(config_.GetSelfInfo().ip());
coord_contact.set_port(config_.GetSelfInfo().port());

if (!cross_peers.empty()) {
  for (const auto& peer : cross_peers) {
    Request prep;
    prep.CopyFrom(*user_request);
    prep.set_type(Request::TYPE_2PC_PREPARE);
    *prep.mutable_client_info() = coord_contact;
    NetChannel channel(peer.ip(), peer.port());
    channel.SendRawMessage(prep);   // direct, short-conn, unsigned
  }
}

// Store request while waiting for votes
{
  std::lock_guard<std::mutex> lk(twopc_mutex_);
  user_request->set_type(Request::TYPE_PRE_PREPARE);   // restore for PBFT
  pending_2pc_requests_[*seq] = std::move(user_request);
  vote_count_[*seq] = 0;
}
return 0;
\end{lstlisting}

\subsection{Participant: voting YES on PREPARE}

The assignment explicitly says no concurrency control is required and no
aborts need to be modelled, so every participant unconditionally returns
\texttt{YES}. The participant also stashes the PREPARE payload, keyed by
request hash, so it can drive its own local PBFT when the corresponding
GLOBAL\_COMMIT arrives.

\begin{lstlisting}[style=cppstyle,caption={\texttt{platform/consensus/ordering/pbft/commitment.cpp} -- \texttt{Process2PCPrepare}: stash PREPARE, send VOTE back to coordinator.}]
int Commitment::Process2PCPrepare(std::unique_ptr<Context> context,
                                  std::unique_ptr<Request> request) {
  uint64_t seq = request->seq();
  int coordinator_id = request->sender_id();
  LOG(ERROR) << "[2PC] Participant " << config_.GetSelfInfo().id()
             << " received PREPARE for seq: " << seq
             << " from coordinator: " << coordinator_id;

  // Cross-shard: stash the PREPARE keyed by hash so the participant can
  // drive its own local PBFT when the matching GLOBAL_COMMIT arrives.
  if (request->has_client_info() && !request->client_info().ip().empty()) {
    std::lock_guard<std::mutex> lk(participant_twopc_mutex_);
    participant_twopc_by_hash_[request->hash()] =
        std::make_unique<Request>(*request);
  }

  // Always vote YES (no aborts, per assignment spec)
  Request vote;
  vote.set_type(Request::TYPE_2PC_VOTE);
  vote.set_seq(seq);
  vote.set_sender_id(config_.GetSelfInfo().id());
  vote.set_current_view(message_manager_->GetCurrentView());
  vote.set_hash(request->hash());
  vote.set_proxy_id(request->proxy_id());

  if (request->has_client_info() && !request->client_info().ip().empty()) {
    const auto& dest = request->client_info();
    NetChannel channel(dest.ip(), dest.port());
    channel.SendRawMessage(vote);
  } else {
    replica_communicator_->SendMessage(vote, coordinator_id);
  }
  return 0;
}
\end{lstlisting}

\subsection{Coordinator: collecting votes and committing}

Once the coordinator has received one YES from every cross-shard peer, it
broadcasts \texttt{2PC GLOBAL\_COMMIT} to those peers and then initiates
the local PBFT round on its own shard. The
\texttt{stored\_request->set\_type(TYPE\_PRE\_PREPARE)} on line~478 transitions
the message back into the standard PBFT pipeline.

\begin{lstlisting}[style=cppstyle,caption={\texttt{platform/consensus/ordering/pbft/commitment.cpp} -- \texttt{Process2PCVote}: count votes, fan out GLOBAL\_COMMIT, start PBFT.}]
int Commitment::Process2PCVote(std::unique_ptr<Context> context,
                               std::unique_ptr<Request> request) {
  uint64_t seq = request->seq();
  int voter_id = request->sender_id();
  LOG(ERROR) << "[2PC] Coordinator received VOTE from replica " << voter_id
             << " for seq: " << seq;

  std::unique_ptr<Request> stored_request;
  bool all_votes_received = false;

  {
    std::lock_guard<std::mutex> lk(twopc_mutex_);
    vote_count_[seq]++;
    const int num_participants =
        config_.GetCrossShardPeers().empty()
            ? static_cast<int>(config_.GetReplicaNum()) - 1
            : static_cast<int>(config_.GetCrossShardPeers().size());

    if (vote_count_[seq] >= num_participants) {
      all_votes_received = true;
      stored_request = std::move(pending_2pc_requests_[seq]);
      pending_2pc_requests_.erase(seq);
      vote_count_.erase(seq);
    }
  }

  if (all_votes_received && stored_request) {
    // === 2PC Phase 2: GLOBAL COMMIT ===
    Request commit_msg;
    commit_msg.set_type(Request::TYPE_2PC_COMMIT);
    commit_msg.set_seq(seq);
    commit_msg.set_sender_id(config_.GetSelfInfo().id());
    commit_msg.set_current_view(message_manager_->GetCurrentView());
    commit_msg.set_hash(stored_request->hash());

    const std::vector<ReplicaInfo> cross_peers =
        config_.GetCrossShardPeers();
    if (!cross_peers.empty()) {
      for (const auto& peer : cross_peers) {
        NetChannel channel(peer.ip(), peer.port());
        channel.SendRawMessage(commit_msg);
      }
    }
    replica_communicator_->BroadCast(commit_msg);

    // === Now run PBFT on the coordinator shard ===
    stored_request->set_type(Request::TYPE_PRE_PREPARE);
    replica_communicator_->BroadCast(*stored_request);
  }
  return 0;
}
\end{lstlisting}

\subsection{Participant: receiving GLOBAL\_COMMIT and starting local PBFT}

When a participant shard leader receives \texttt{2PC GLOBAL\_COMMIT}, it
retrieves the stashed PREPARE, assigns a fresh local sequence number, and
broadcasts \texttt{PRE\_PREPARE} to its own four replicas. Each shard
therefore runs an independent PBFT round for the same transaction (the
``every shard executes'' option from the assignment spec).

\begin{lstlisting}[style=cppstyle,caption={\texttt{platform/consensus/ordering/pbft/commitment.cpp} -- \texttt{Process2PCCommit}: drive local PBFT from the stashed PREPARE.}]
int Commitment::Process2PCCommit(std::unique_ptr<Context> context,
                                 std::unique_ptr<Request> request) {
  uint64_t seq = request->seq();
  LOG(ERROR) << "[2PC] Participant " << config_.GetSelfInfo().id()
             << " received GLOBAL COMMIT for seq: " << seq;

  std::unique_ptr<Request> stashed;
  const std::string commit_hash = request->hash();
  {
    std::lock_guard<std::mutex> lk(participant_twopc_mutex_);
    auto it = participant_twopc_by_hash_.find(commit_hash);
    if (it != participant_twopc_by_hash_.end()) {
      stashed = std::move(it->second);
      participant_twopc_by_hash_.erase(it);
    }
  }
  if (!stashed) {
    return 0;   // local backups have no stash; coordinator drives them.
  }

  auto local_seq = message_manager_->AssignNextSeq();
  if (!local_seq.ok()) {
    std::lock_guard<std::mutex> lk(participant_twopc_mutex_);
    participant_twopc_by_hash_[commit_hash] = std::move(stashed);
    return -2;
  }
  stashed->set_seq(*local_seq);
  stashed->set_current_view(message_manager_->GetCurrentView());
  stashed->set_sender_id(config_.GetSelfInfo().id());
  stashed->set_primary_id(message_manager_->GetCurrentPrimary());
  stashed->set_type(Request::TYPE_PRE_PREPARE);
  replica_communicator_->BroadCast(*stashed);
  return 0;
}
\end{lstlisting}

\subsection{Participant: accepting cross-shard PRE\_PREPAREs}

When a participant shard broadcasts the stashed transaction's
\texttt{PRE\_PREPARE} to its own backups (Section~3.7), the backups
verify \texttt{request->data\_signature()} in
\texttt{Commitment::ProcessProposeMsg}. The signature was produced by
the coordinator shard's leader, whose public key the participant
shard's verifier does not hold -- shards exchange keys only within
their own region. The signature check therefore distinguishes
local from foreign signers: if the signer is a member of the local
shard the check is enforced as usual, but a signer from another shard
is accepted unconditionally (the assignment explicitly waives
concurrency control and aborts on cross-shard work). This is what
allows a \texttt{SET} routed to one coordinator to be readable from
any of the four shards in a later \texttt{GET}.

\begin{lstlisting}[style=cppstyle,caption={\texttt{platform/consensus/ordering/pbft/commitment.cpp} -- accept cross-shard \texttt{PRE\_PREPARE} (foreign signer).}]
bool valid =
    verifier_->VerifyMessage(request->data(), request->data_signature());
if (!valid) {
  // Cross-shard PRE_PREPARE: signed by another shard's leader whose
  // key this verifier does not know. Accept iff signer is foreign.
  int64_t signer_id = request->data_signature().node_id();
  bool signer_is_local = false;
  for (const auto& r : config_.GetReplicaInfos()) {
    if (r.id() == signer_id) { signer_is_local = true; break; }
  }
  if (signer_is_local) {
    LOG(ERROR) << "request is not valid: "
               << request->data_signature().DebugString();
    return -2;
  }
  // foreign signer -> trust per spec ("no concurrency control / no aborts")
}
\end{lstlisting}

\subsection{Enabling the client-request drain on shard leaders}

\texttt{ResponseManager} runs a background \texttt{BatchProposeMsg}
thread that drains the per-leader \texttt{batch\_queue\_} and forwards
batched user requests to the primary as \texttt{TYPE\_NEW\_TXNS}. The
sharded deployment has no separate proxy node -- the shard leaders
themselves receive client transactions -- so the drain thread runs on
any node with a non-empty \texttt{crossShardPeer} list (the shard
leaders, by construction). This is what wires client requests into
the 2PC + PBFT pipeline on a shard leader.

\begin{lstlisting}[style=cppstyle,caption={\texttt{platform/consensus/ordering/pbft/response\_manager.cpp} -- drain thread also runs on sharded leaders.}]
if (config_.GetPublicKeyCertificateInfo()
            .public_key()
            .public_key_info()
            .type() == CertificateKeyInfo::CLIENT ||
    config_.IsTestMode() ||
    !config_.GetCrossShardPeers().empty()) {  // sharded leaders
  user_req_thread_ = std::thread(&ResponseManager::BatchProposeMsg, this);
}
\end{lstlisting}

\section{Build and run}

The whole system runs from a single Docker image, \texttt{dom-sharded}, built
from \texttt{Docker/Dockerfile\_mac} on an Apple Silicon host (or
\texttt{Docker/Dockerfile} with \texttt{--platform linux/amd64} on x86):

\begin{lstlisting}[style=shstyle,caption={Build and run.}]
# Build
docker build -t dom-sharded -f Docker/Dockerfile_mac .

# Run (entrypoint brings up all 16 nodes and tails)
docker run -it --name dom-sharded dom-sharded

# In a second shell: drive a SET / GET
docker exec dom-sharded \
  /app/bazel-bin/service/tools/kv/api_tools/kv_service_tools \
  /app/scripts/deploy/config/sharded/client.shard_leaders.config \
  set k1 v1

docker exec dom-sharded \
  /app/bazel-bin/service/tools/kv/api_tools/kv_service_tools \
  /app/scripts/deploy/config/sharded/client.shard_leaders.config \
  get k1

# Long-lived multi-threaded benchmark (rotates across all 4 leaders).
# Modes: --ops N (fixed) or --duration SEC (steady-state flood).
docker exec dom-sharded \
  /app/bazel-bin/service/tools/kv/api_tools/kv_benchmark \
  --config /app/scripts/deploy/config/sharded/client.shard_leaders.config \
  --duration 60 --concurrency 8 --prefix mt_ --verify 0
\end{lstlisting}

\section{Evaluation}

\subsection{Machine configuration}

\begin{tabular}{ll}
  \toprule
  Host             & Apple MacBook (M-series, Apple Silicon) \\
  Container        & \texttt{dom-sharded}, \texttt{arm64v8/ubuntu:20.04} \\
  Cluster topology & 4 shards $\times$ 4 replicas = 16 processes, all on \texttt{127.0.0.1} \\
  Compiler         & Clang via Bazel 6.0.0 (Bazelisk) \\
  \bottomrule
\end{tabular}

\medskip
\noindent\textit{Replace the host line above with the exact CPU model, core count,
RAM, OS version, and Docker Desktop version from} \texttt{system\_profiler
SPHardwareDataType} \textit{and} \texttt{docker version} \textit{before
submitting.}

\subsection{Correctness trace (single transaction)}

The following lines from the shard-1 leader log
(\texttt{kv\_1.log}) show one transaction proceeding through both 2PC phases
and into local PBFT:

\begin{lstlisting}[style=shstyle,caption={Annotated coordinator log for one transaction.}]
[2PC] Coordinator starting 2PC for seq: 1
[2PC] sending PREPARE seq=1 to peer id=5  ip=127.0.0.1 port=18011
[2PC] sending PREPARE seq=1 to peer id=9  ip=127.0.0.1 port=18021
[2PC] sending PREPARE seq=1 to peer id=13 ip=127.0.0.1 port=18031
[2PC] Coordinator received VOTE from replica 5  for seq: 1
[2PC] Coordinator received VOTE from replica 9  for seq: 1
[2PC] Coordinator received VOTE from replica 13 for seq: 1
[2PC] Vote count for seq 1: 3/3
[2PC] All votes received for seq: 1. Broadcasting GLOBAL COMMIT.
[2PC] GLOBAL_COMMIT -> peer id=5  / 9 / 13
[2PC] 2PC complete for seq: 1. Starting PBFT consensus.
\end{lstlisting}

\subsection{Throughput and latency}

A long--lived multi--threaded C++ benchmark driver
(\texttt{service/tools/kv/api\_tools/kv\_benchmark.cpp}) supports two
modes: a fixed \texttt{--ops N} count and a steady--state
\texttt{--duration SEC} flood, both with \texttt{--concurrency N}
client threads. Each thread holds one \texttt{KVClient} per shard
leader and rotates among them, so coordinator duty fans out across
all four shards in parallel. A host--side wrapper
(\texttt{DomNotes/run\_benchmark.sh}) builds the binary inside the
\texttt{dom-sharded} container.

\paragraph{Per-shard server-side throughput.} Per-shard
throughput is parsed from each primary's \texttt{stats.cpp:424}
monitor lines (which emit a \texttt{total request:N txn:N time:T}
record every 5\,s). Average TPS averages the active 5\,s windows;
max TPS is the peak window; latency is the average of
\texttt{stats.cpp:474} \texttt{req client latency} lines.

\begin{tabular}{lcccc}
  \toprule
  & \textbf{Avg TPS} & \textbf{Max TPS} & \textbf{Avg latency (s)} \\
  \midrule
  Shard 1 (\texttt{kv\_1})  & 82.8 & 110 & $7.9 \times 10^{-5}$ \\
  Shard 2 (\texttt{kv\_5})  & 82.0 & 110 & $8.0 \times 10^{-5}$ \\
  Shard 3 (\texttt{kv\_9})  & 81.4 & 110 & $8.0 \times 10^{-5}$ \\
  Shard 4 (\texttt{kv\_13}) & 80.6 & 110 & $7.9 \times 10^{-5}$ \\
  \midrule
  \textbf{Mean per shard}   & \textbf{81.7} & \textbf{110} & $\mathbf{\sim 7.9 \times 10^{-5}}$ \\
  \textbf{Aggregate cluster}& \textbf{326.8} & \textbf{$\sim$440} & --- \\
  \bottomrule
\end{tabular}

\medskip
\paragraph{Per-shard client-observed throughput.} For
completeness, the same 60\,s 8-thread runs measured client-side give:

\begin{tabular}{lcccc}
  \toprule
   & Trial 1 & Trial 2 & Trial 3 & \textbf{Mean} \\
  \midrule
  Ops succeeded                & 5{,}159 & 6{,}053 & 5{,}510 & 5{,}574 \\
  Throughput (txn/s)           & 85.3    & 100.1   & 91.2    & \textbf{92.2} \\
  Avg latency (ms / op)        & 11.72   & 9.99    & 10.97   & \textbf{10.89} \\
  \bottomrule
\end{tabular}

\medskip
\begin{tabular}{lr}
  \toprule
  \multicolumn{2}{l}{\textbf{Durability check: 20 \texttt{SET}s then 20 \texttt{GET}s, single thread}} \\
  \midrule
  Keys written (round--robin across 4 shards)           & 20 \\
  Keys read back (round--robin across 4 shards)         & 20 \\
  Hits / misses                                         & \textbf{20 / 0} \\
  \texttt{GET} average latency (ms / op)                & 157.7 \\
  \bottomrule
\end{tabular}

\medskip
\noindent\textbf{Reading the numbers.}
\textit{Two measurements, two stories.} The server-side per-shard TPS
(81.7 average, 326.8 aggregate) and the client-observed end-to-end
TPS (92.2) measure different things. The server-side number counts
client requests arriving at each primary divided by the 5\,s monitor
window, then averaged across active windows. The client-observed
number is wall-clock ops/second around the \texttt{Set()} loop.
The client's \texttt{Set()} call is fire-and-forget at the TCP layer
(\texttt{transaction\_constructor.cpp}'s single-arg
\texttt{SendRequest} passes \texttt{wait\_response=false}) and each
call opens a fresh \texttt{NetChannel}, so the kernel's local-loopback
3-way handshake under 8 contending threads is the practical ceiling
for client-observed throughput on a single host.

The four shards' server-side TPS are within $\pm 1$\% of each other,
confirming that the proxy's round-robin (Section~3.3) genuinely
spreads coordinator duty equally. Max windows reach $\sim$110\,TPS,
i.e., the per-shard ceiling is $\sim$1.3$\times$ the steady-state
average -- a healthy headroom indicator. Server-internal latency
averages $7.9 \times 10^{-5}$\,s ($\sim$79\,$\mu$s), which is one
stage of the pipeline (the time between request arrival and PBFT
commit), not the end-to-end RTT. The end-to-end client-observed
latency of $\sim$10.9\,ms includes TCP setup + serialization +
2PC + PBFT + reply.

\textit{Durability.} The 20--op SET+GET pass writes a key against one
shard's coordinator, then reads it back via a round--robin'd \texttt{GET}
that may land on \emph{any} of the four shards. All 20 reads hit,
meaning every cross--shard \texttt{GLOBAL\_COMMIT} did successfully
drive local PBFT on every participant shard. Internally,
\texttt{ResponseManager}'s \texttt{BatchProposeMsg} groups client
requests into batches, so roughly 5,700 ops per trial produce
$\sim$700--900 server--side batched commits per trial (summed across
the four shards) rather than one round per op. Every round carries
a full cross--shard 2PC (PREPARE $\rightarrow$ 3$\times$VOTE
$\rightarrow$ GLOBAL\_COMMIT) followed by intra--shard PBFT on
\emph{all four shards}.

\subsection{Notes on the throughput methodology}

The fire--and--forget \texttt{Set()} path limits the comparability of
the client--observed throughput against systems that report
commit--acknowledged TPS. We therefore report both: client--observed
TPS for honest end--to--end RTT context, and server--side TPS for
per--shard steady--state measurement. A small helper script
\texttt{DomNotes/parse\_server\_stats.py} extracts the per--shard
table from \texttt{kv\_\{1,5,9,13\}.log} files in
\texttt{config\_out\_sharded/logs/}; it computes Avg TPS by averaging
\texttt{total request:N / 5\,s} across active monitor windows and Max
TPS by taking the peak window.

The coordinator's \texttt{Process2PCVote} tracks votes using a
\texttt{std::set<int64\_t> twopc\_voters\_} keyed by sequence number,
which records distinct voter IDs. Quorum is reached only when the
set size meets the participant count, so duplicate votes from the same
participant cannot prematurely trip the ``all votes received''
threshold under concurrent client load.

\section{Implementation by file}

\begin{tabular}{lp{0.6\linewidth}}
  \toprule
  \textbf{File} & \textbf{Role} \\
  \midrule
  \texttt{platform/proto/resdb.proto}
    & Declares \texttt{TYPE\_2PC\_PREPARE/VOTE/COMMIT} (values 22--24) as request types alongside the existing PBFT types. \\
  \texttt{platform/proto/replica\_info.proto}
    & Declares the \texttt{cross\_shard\_peer} list and \texttt{multi\_shard\_client\_round\_robin} flag on \texttt{ResConfigData}. \\
  \texttt{platform/config/resdb\_config.\{h,cpp\}}
    & Exposes \texttt{GetCrossShardPeers()} and \texttt{MultiShardClientRoundRobin()} accessors for the runtime to read the new fields. \\
  \texttt{interface/rdbc/transaction\_constructor.\{h,cpp\}}
    & \texttt{PickDestReplica()} performs round-robin across shard leaders when configured. The rotor counter is \texttt{static} so all \texttt{KVClient} instances in one process share a single round-robin pointer. \\
  \texttt{platform/consensus/ordering/pbft/commitment.\{h,cpp\}}
    & Implements \texttt{Process2PCPrepare/Vote/Commit}, the \texttt{pending\_2pc\_requests\_} and \texttt{participant\_twopc\_by\_hash\_} state, the coordinator's PREPARE/GLOBAL\_COMMIT broadcasts, and the participant's local-PBFT kick-off. \texttt{ProcessProposeMsg} accepts \texttt{PRE\_PREPARE}s signed by a foreign-shard leader (and enforces the signature for local-shard signers). Vote tracking uses \texttt{std::set<int64\_t> twopc\_voters\_} so duplicate votes do not trip early quorum. Cross-shard sends use direct \texttt{NetChannel} writes. \\
  \texttt{platform/consensus/ordering/pbft/consensus\_manager\_pbft.cpp}
    & Dispatches the three 2PC request types to their handlers in \texttt{Commitment}. \\
  \texttt{platform/consensus/ordering/pbft/response\_manager.cpp}
    & \texttt{BatchProposeMsg} drains the per-leader queue and forwards batched user requests; runs on any node that has \texttt{crossShardPeer} entries (i.e., the shard leaders). \\
  \texttt{scripts/deploy/config/sharded/}
    & Four \texttt{shard\{1..4\}.server.config} files (each shard plus its 3 cross-shard peers) and a \texttt{client.shard\_leaders.config} for the round-robin client. \\
  \texttt{entrypoint\_sharded.sh}, \\
  \texttt{scripts/deploy/script/start\_sharded\_kv\_cluster.sh}, \\
  \texttt{scripts/deploy/script/generate\_sharded\_configs.sh}
    & Generate certificates, bring up the 16-node cluster, and serve as the container entrypoint. \\
  \texttt{Docker/Dockerfile\_mac}
    & Builds the \texttt{dom-sharded} image natively on Apple Silicon, with sharded scripts marked executable and \texttt{entrypoint\_sharded.sh} as the default \texttt{ENTRYPOINT}. \\
  \texttt{service/tools/kv/api\_tools/kv\_benchmark.cpp},\\
  \texttt{service/tools/kv/api\_tools/BUILD},\\
  \texttt{DomNotes/run\_benchmark.sh},\\
  \texttt{DomNotes/parse\_server\_stats.py}
    & Long-lived multi-threaded C++ benchmark driver (\texttt{--ops} and \texttt{--duration} modes, \texttt{--concurrency N}, one \texttt{KVClient} per leader); a host-side runner; and a Python parser that extracts per-shard Avg/Max TPS and avg latency from each primary's \texttt{stats.cpp} monitor lines. \\
  \bottomrule
\end{tabular}

\section{Conclusion}

Every required architectural element of the assignment is implemented and
verified end to end. The cluster runs 4 shards of 4 replicas, the proxy
alternates client requests across the four shard leaders, every
transaction triggers a cross-shard 2PC PREPARE--VOTE--GLOBAL\_COMMIT
exchange, and the GLOBAL\_COMMIT then drives an independent PBFT round on
each shard. A long--lived, multi--threaded client distributes coordinator duty
evenly across all four shard leaders: averaged over three 60--second
trials at concurrency 8, each shard's server--side average throughput
sits within $\pm 1$\% of the others
($\sim$\textbf{82\,TPS} per shard, \textbf{$\sim$327\,TPS} aggregate
cluster, peak windows $\sim$110\,TPS). A 20--operation
\texttt{SET}+\texttt{GET} pass returns \textbf{20 of 20} keys
regardless of which shard the \texttt{GET} is round--robin'd to --
direct evidence that the cross--shard \texttt{GLOBAL\_COMMIT}
successfully drives local PBFT on every participant shard.
Server--internal latency is $\sim$\textbf{79\,$\mu$s} per request;
end--to--end client--observed latency is $\sim$\textbf{10.9\,ms},
dominated by per--call TCP setup on the client side.

\end{document}
